% GENERATED FILE. Edit canonical lessons, then run npm run generate:books.
% canonical-source-hash: fnv1a64:0b08d317aa3a378f
% canonical-lessons: LA-C391-etiamsi, LA-C391-tametsi, LA-C391-alioqui, LA-C391-secus, LA-C391-sive

\chapter{Even If and Otherwise}
\label{ch:la-even-if-and-otherwise}
\hlchaptermodality{391}

\begin{chapteropening}
\textbf{By the end of this chapter:} \emph{I can say even if and although, otherwise and differently, and whether … or.}

Complete the last lesson of chapter 391: I can say even if and although, otherwise and differently, and whether … or.
\end{chapteropening}

\section[etiamsī]{etiamsī — even if}

\label{lesson:LA-C391-etiamsi}



\begin{quote}
Before the new one: say the Latin for only, then the Latin for however much.
\end{quote}

\subsection*{You'll want to know: etiamsī}
\textbf{etiamsī} — \textquotedblleft{}even if\textquotedblright{}.

\begin{quote}
Etiamsī pluet, ambulābimus.
\end{quote}

\begin{quote}
\emph{Even if it rains, we shall walk.}
\end{quote}

\emph{Etiamsī} = \emph{etiam} + \emph{sī}: even if. It takes the mood a \emph{sī} clause would take.

\begin{cousinweb}
\emph{etiamsī} is built from \emph{etiam} \textquotedblleft{}even, also\textquotedblright{} and \emph{sī} \textquotedblleft{}if\textquotedblright{} (Lewis \& Short), and is often written as two words.
\end{cousinweb}

\begin{grammarlens}[title={What you've built}]
One more word for telling, explaining and describing at B1.
\end{grammarlens}

\subsection*{Your turn}
\begin{itemize}
\raggedright
  \item \emph{Say it:} \emph{etiamsī}
  \item \emph{Say it:} \emph{etiamsī}, once more
  \item \emph{Say it:} a sentence of your own with \emph{etiamsī}
  \item \emph{Recall:} say the Latin for only, then the Latin for however much, then say \emph{etiamsī} again
\end{itemize}

\begin{tcolorbox}[breakable,colback=teal!4,colframe=teal!35!black,title={Before you move on}]
What does etiamsī mean? (\textquotedblleft{}Even if\textquotedblright{}.) Say it once more.
\end{tcolorbox}
\section[tametsī]{tametsī — although, even though}

\label{lesson:LA-C391-tametsi}



\begin{quote}
Before the new one: say the Latin for however much, then the Latin for even if.
\end{quote}

\subsection*{You'll want to know: tametsī}
\textbf{tametsī} — \textquotedblleft{}although, even though\textquotedblright{}.

\begin{quote}
Tametsī fessus erat, labōrāvit.
\end{quote}

\begin{quote}
\emph{Although he was tired, he worked.}
\end{quote}

\emph{Tametsī} = \emph{tamen} + \emph{etsī}, and like \emph{etsī} it mostly takes the indicative.

\begin{cousinweb}
\emph{tametsī} is built from \emph{tamen} \textquotedblleft{}yet\textquotedblright{} and \emph{etsī} \textquotedblleft{}even if\textquotedblright{} (Lewis \& Short), with the end of \emph{tamen} worn away.
\end{cousinweb}

\begin{grammarlens}[title={What you've built}]
One more word for telling, explaining and describing at B1.
\end{grammarlens}

\subsection*{Your turn}
\begin{itemize}
\raggedright
  \item \emph{Say it:} \emph{tametsī}
  \item \emph{Say it:} \emph{tametsī}, once more
  \item \emph{Say it:} a sentence of your own with \emph{tametsī}
  \item \emph{Recall:} say the Latin for however much, then the Latin for even if, then say \emph{tametsī} again
\end{itemize}

\begin{tcolorbox}[breakable,colback=teal!4,colframe=teal!35!black,title={Before you move on}]
What does tametsī mean? (\textquotedblleft{}Although, even though\textquotedblright{}.) Say it once more.
\end{tcolorbox}
\section[aliōquī]{aliōquī — otherwise; in other respects}

\label{lesson:LA-C391-alioqui}



\begin{quote}
Before the new one: say the Latin for even if, then the Latin for although.
\end{quote}

\subsection*{You'll want to know: aliōquī}
\textbf{aliōquī} — \textquotedblleft{}otherwise; in other respects\textquotedblright{}.

\begin{quote}
Festīnāre dēbēmus; aliōquī nāvem āmittēmus.
\end{quote}

\begin{quote}
\emph{We must hurry; otherwise we shall miss the ship.}
\end{quote}

\emph{Aliōquī} (also \emph{aliōquīn}): otherwise, or else; also, in other respects.

\begin{cousinweb}
\emph{aliōquī} is built from a form of \emph{alius} \textquotedblleft{}other\textquotedblright{} and the old adverb \emph{quī} \textquotedblleft{}in some way\textquotedblright{} (EtymDB; Lewis \& Short): in another way.
\end{cousinweb}

\begin{grammarlens}[title={What you've built}]
One more word for telling, explaining and describing at B1.
\end{grammarlens}

\subsection*{Your turn}
\begin{itemize}
\raggedright
  \item \emph{Say it:} \emph{aliōquī}
  \item \emph{Say it:} \emph{aliōquī}, once more
  \item \emph{Say it:} a sentence of your own with \emph{aliōquī}
  \item \emph{Recall:} say the Latin for even if, then the Latin for although, then say \emph{aliōquī} again
\end{itemize}

\begin{tcolorbox}[breakable,colback=teal!4,colframe=teal!35!black,title={Before you move on}]
What does aliōquī mean? (\textquotedblleft{}Otherwise; in other respects\textquotedblright{}.) Say it once more.
\end{tcolorbox}
\section[secus]{secus — differently, otherwise}

\label{lesson:LA-C391-secus}



\begin{quote}
Before the new one: say the Latin for although, then the Latin for otherwise.
\end{quote}

\subsection*{You'll want to know: secus}
\textbf{secus} — \textquotedblleft{}differently, otherwise\textquotedblright{}.

\begin{quote}
Rēs secus accidit atque spērābāmus.
\end{quote}

\begin{quote}
\emph{Things turned out otherwise than we hoped.}
\end{quote}

\emph{Secus} is an adverb: \emph{nōn secus ac}, not otherwise than, in the same way as; \emph{haud secus}, in no other way.

\begin{cousinweb}
\emph{secus} comes from the root of \emph{sequor} \textquotedblleft{}to follow\textquotedblright{} (EtymDB; Lewis \& Short): coming after, so less well, otherwise.
\end{cousinweb}

\begin{grammarlens}[title={What you've built}]
One more word for telling, explaining and describing at B1.
\end{grammarlens}

\subsection*{Your turn}
\begin{itemize}
\raggedright
  \item \emph{Say it:} \emph{secus}
  \item \emph{Say it:} \emph{secus}, once more
  \item \emph{Say it:} a sentence of your own with \emph{secus}
  \item \emph{Recall:} say the Latin for although, then the Latin for otherwise, then say \emph{secus} again
\end{itemize}

\begin{tcolorbox}[breakable,colback=teal!4,colframe=teal!35!black,title={Before you move on}]
What does secus mean? (\textquotedblleft{}Differently, otherwise\textquotedblright{}.) Say it once more.
\end{tcolorbox}
\section[sīve]{sīve — or if; whether … or}

\label{lesson:LA-C391-sive}



\begin{quote}
Before the new one: say the Latin for otherwise, then the Latin for differently.
\end{quote}

\subsection*{You'll want to know: sīve}
\textbf{sīve} — \textquotedblleft{}or if; whether … or\textquotedblright{}.

\begin{quote}
Sīve pluit sīve sōl lūcet, ambulō.
\end{quote}

\begin{quote}
\emph{Whether it rains or the sun shines, I go for a walk.}
\end{quote}

\emph{Sīve} (or \emph{seu}) = \emph{sī} + \emph{-ve}, or: \emph{sīve … sīve}, whether … or.

\begin{cousinweb}
\emph{sīve} is built from \emph{sī} \textquotedblleft{}if\textquotedblright{} and the particle \emph{-ve} \textquotedblleft{}or\textquotedblright{} (Lewis \& Short).
\end{cousinweb}

\begin{grammarlens}[title={What you've built}]
One more word for telling, explaining and describing at B1.
\end{grammarlens}

\subsection*{Your turn}
\begin{itemize}
\raggedright
  \item \emph{Say it:} \emph{sīve}
  \item \emph{Say it:} \emph{sīve}, once more
  \item \emph{Say it:} a sentence of your own with \emph{sīve}
  \item \emph{Recall:} say the Latin for otherwise, then the Latin for differently, then say \emph{sīve} again
\end{itemize}

\begin{tcolorbox}[breakable,colback=teal!4,colframe=teal!35!black,title={Before you move on}]
What does sīve mean? (\textquotedblleft{}Or if; whether … or\textquotedblright{}.) Say it once more.
\end{tcolorbox}
